%==============================================================================
%  WEEKLY AI BRIEF -- issue #10, week ending 27 October 2026
%  Engine: tectonic
%==============================================================================
\documentclass[11pt,a4paper]{article}

%------------------------------ packages -------------------------------------
\usepackage[T1]{fontenc}
\usepackage[utf8]{inputenc}
\usepackage{lmodern}
\usepackage[margin=2.2cm,top=2.0cm,bottom=2.3cm]{geometry}
\usepackage{microtype}
\usepackage{graphicx}
\usepackage{xcolor}
\usepackage{titlesec}
\usepackage{enumitem}
\usepackage{booktabs}
\usepackage{tabularx}
\usepackage{fancyhdr}
\usepackage{lastpage}
\usepackage{tcolorbox}
\usepackage{xurl}
\usepackage{needspace}
\usepackage[hidelinks]{hyperref}

\tcbuselibrary{skins,breakable}

% Logos live in template/ and are shared by every issue -- never copy them into a
% dated folder. This path list resolves them from <YYYY-MM-DD>/ or from template/.
\graphicspath{{../../../template/}{./}{template/}}

%------------------------------ palette --------------------------------------
\definecolor{brandink}{HTML}{103868}    % KAUST navy -- headings, rules
\definecolor{brandaccent}{HTML}{0B7480} % KAUST teal (darkened to 5.6:1 on white)
\definecolor{brandsoft}{HTML}{EAF3F5}   % box fills
\definecolor{brandgrey}{HTML}{5A6672}   % meta text

\hypersetup{colorlinks=true,linkcolor=brandaccent,urlcolor=brandaccent,
            citecolor=brandaccent,pdfborder={0 0 0}}

%------------------------------ metadata -------------------------------------
% Issue number: one more than the highest \briefnumber in the existing dated folders.
% A REBUILD of a window that already has a folder keeps that folder's number -- the
% number identifies the window, not the run.
\newcommand{\briefnumber}{10}
% The window is seven days ending on \briefdate, INCLUSIVE -- so the span is 6 days
% wide (19--25, not 18--25). Check both endpoints against `date +%F` before shipping.
\newcommand{\briefweek}{21--27 October 2026}
\newcommand{\briefdate}{27 October 2026}            % date of publication
\newcommand{\briefauthor}{Ali Habibullah}          % prepared by
\newcommand{\briefaffil}{KAUST}                    % affiliation
\newcommand{\briefscope}{Papers, model releases, benchmarks and policy from the past seven days}

%------------------------------ section styling ------------------------------
\titleformat{\section}
  {\normalfont\sffamily\large\bfseries\color{brandink}}{\thesection.}{0.6em}{}
  [\vspace{-0.55em}{\color{brandaccent}\rule{\textwidth}{0.9pt}}]
\titlespacing*{\section}{0pt}{1.5ex plus .3ex}{0.9ex}

\titleformat{\subsection}[runin]
  {\normalfont\sffamily\bfseries\color{brandink}}{}{0pt}{}[\;---]
\titlespacing*{\subsection}{0pt}{1.25ex plus .2ex}{0.5em}

\setlist[itemize]{leftmargin=1.15em,itemsep=0.25em,topsep=0.35em,parsep=0pt}
\setlength{\parindent}{0pt}
\setlength{\parskip}{0.45em}
\setlength{\emergencystretch}{1.5em}

%------------------------------ header / footer ------------------------------
\pagestyle{fancy}
\fancyhf{}
\renewcommand{\headrulewidth}{0pt}
\renewcommand{\footrulewidth}{0.4pt}
\fancyfoot[L]{\footnotesize\color{brandgrey}Weekly AI Brief \#\briefnumber\ \textbullet\ \briefweek}
\fancyfoot[R]{\footnotesize\color{brandgrey}Page \thepage\ of \pageref{LastPage}}

%------------------------------ content macros -------------------------------
% Executive-summary box.  \tldr{ \item ... \item ... }
\newcommand{\tldr}[1]{%
  \begin{tcolorbox}[breakable,enhanced,colback=brandsoft,colframe=brandaccent,
      boxrule=0.8pt,left=10pt,right=10pt,top=8pt,bottom=8pt,arc=2pt,
      title={\sffamily\bfseries The week in five lines},fonttitle=\color{white}]
    \begin{itemize}[leftmargin=1.1em,itemsep=0.3em,topsep=0pt]#1\end{itemize}
  \end{tcolorbox}}

% One item of news.  \entry{Headline}{Who / venue}{URL}{source-number}
\newcommand{\entry}[4]{%
  \par\Needspace{7\baselineskip}\vspace{0.7em}%
  {\sffamily\bfseries\color{brandink}#1}\par
  {\footnotesize\color{brandgrey}#2 \textbullet\ \href{#3}{link} \textbullet\ \srcref{#4}}\par
  \vspace{0.15em}}

% "Why it matters" line -- one sentence, no hedging.
\newcommand{\why}[1]{\par{\small\textbf{Why it matters.} #1}\par}

% Numbered pointer into SOURCES.md / the source table.
\newcommand{\srcref}[1]{\textcolor{brandaccent}{[S#1]}}

% Numbers worth pulling out of a paragraph.
\newcommand{\stat}[2]{{\sffamily\bfseries\color{brandink}#1}~{\footnotesize\color{brandgrey}#2}}

%==============================================================================
\begin{document}

%------------------------------ title block ----------------------------------
\noindent
\begin{minipage}[c]{0.48\textwidth}
  \raggedright\includegraphics[height=1.4cm]{KAUST_Academy_Logo.png}
\end{minipage}\hfill
\begin{minipage}[c]{0.48\textwidth}
  \raggedleft\includegraphics[height=1.5cm]{KAUST_Logo.png}
\end{minipage}
\vspace{0.7em}

{\sffamily
 {\color{brandgrey}\footnotesize ISSUE \#\briefnumber\ \textbullet\ \briefdate}\par
 \vspace{0.15em}
 {\LARGE\bfseries\color{brandink}Weekly AI Brief}\par
 \vspace{0.25em}
 {\large\color{brandaccent}\briefweek}\par
 \vspace{0.35em}
 {\footnotesize\color{brandgrey}\briefscope\\
  Prepared by \briefauthor, \briefaffil}\par}
\vspace{0.3em}
{\color{brandink}\rule{\textwidth}{1.4pt}}
\vspace{0.6em}

%------------------------------ 1. executive summary -------------------------
\tldr{
  \item \textbf{Anthropic filed its S-1 to list on Nasdaq as ``ANT''.} The filing shows revenue of \$16.3B for the first half of 2026, up from \$4.6B for all of 2025, and a \$693M operating loss. \srcref{15}\,\srcref{16}
  \item \textbf{Moonshot's Kimi K3.1 Pro scores 55 on Artificial Analysis's index,} fourth overall and the highest-scoring model whose maker has promised open weights. Moonshot says the weights will be released by 4 November. \srcref{1}\,\srcref{2}
  \item \textbf{OpenAI's GPT-6.1 Terra costs \$1/\$5 per million tokens} and scores 48 on the same index, 4 points below GPT-6.1 Sol at about half the price. \srcref{3}\,\srcref{5}
  \item \textbf{OpenAI disclosed that an internal model reached the live internet through a DNS misconfiguration on 20 September.} It now trains frontier models with no internet access. \srcref{18}
  \item \textbf{Meta set out a plan to raise about \$95B in new equity for AI spending:} \$35B in underwritten offerings now and a planned \$60B at-the-market programme later. \srcref{17}
}

%------------------------------ 2. models & releases -------------------------
\section{Model and product releases}

\entry{Kimi K3.1 Pro: fourth on Artificial Analysis, with weights promised by 4 November}{Moonshot AI \textbullet\ API now, weights promised \textbullet\ 23 October}{https://www.kimi.com/en/blog/kimi-k3-1-pro}{1}
Kimi K3.1 Pro is a mixture-of-experts model with 2.8T total and 104B active parameters. It has a 1M-token context and is
available now through the API at \$2/\$10 per million input/output tokens, a third less than Kimi~K3. \srcref{2} Moonshot writes that
``the full model weights will be released by Nov 4''. It does not name a licence. Moonshot reports 65.7\% on Terminal-Bench~4.0
using the Claude Code harness, just behind Claude Opus~5.5 at 66.4\%. \srcref{1} The independent results agree with this. Artificial Analysis scores the model 55 on its index, up from 44 for
K3. That puts it fourth, behind Opus~5.5, Sonnet~5.5 and Muse Wonder Full. AA also measured 67\% on Terminal-Bench~4.0, the highest it has recorded for any model. \srcref{2}
Vals, using its own agent harness, measures 63.64\%, third behind Opus~5.5 and Sonnet~5.5. \srcref{13}
\why{If the weights arrive as promised, the best open-weight model will score within 3 points of the closed frontier on AA's index, against a 12-point gap in last week's issue.}

\entry{GPT-6.1 Terra: OpenAI's cheaper tier, at half the price of Sol}{OpenAI \textbullet\ API and ChatGPT \textbullet\ 26 October}{https://openai.com/index/introducing-gpt-6-1-terra/}{3}
\texttt{gpt-6.1-terra} costs \$1 per million input tokens, \$0.05 for cached input and \$5 for output. OpenAI says that is 50\% below
GPT-5.6 Terra and GPT-6.1 Sol. OpenAI reports 50.3\% on Terminal-Bench~4.0 at maximum effort, up from 21.5\% for GPT-5.6 Terra. It also rates the model at the
``Critical'' cybersecurity threshold of its Preparedness Framework. \srcref{3} Artificial Analysis scores it 48 at maximum effort,
6 points above GPT-5.6 Terra and 4 below Sol. A full run of the index costs \$0.40, against \$1.40 for GPT-5.6 Terra. AA also finds that the model guesses instead of declining on 38\% of questions it cannot answer, down from 88\% for GPT-5.6 Terra. \srcref{4}\,\srcref{5} Vals scores it 60.17\% on its index,
at \$2.92 per test. \srcref{11}
\why{Teams running GPT-5.6 Terra get a large accuracy gain at half the price. The Critical cyber rating means OpenAI's stricter safeguards also apply to this cheaper tier.}

%------------------------------ 3. research ----------------------------------
\section{Research worth reading}

\entry{Inherited, not invented: reasoning models' ``wait'' comes from data, not RL}{Kou et al.\ (incl.\ Yike Guo), HKUST \textbullet\ arXiv:2610.30882 \textbullet\ 22 October}{https://arxiv.org/abs/2610.30882}{6}
In 18 RL runs on base models that rarely say ``wait'', with no supervised reasoning traces first, MATH-500 accuracy rises by a median 42.9 points. In 17 of those runs the word stays below 0.05 uses per
1,000 tokens. Open-Reasoner-Zero-32B, for example, goes from 10.8\% to 91.4\% while writing ``wait'' four times in about 2.4M tokens. The habit instead arrives with
mid-training or fine-tuning on reasoning traces: 13 of 14 fine-tuned models acquire it. A practical warning comes with this. An LLM judge's ``revision'' rate
fell from 26.0\% to 12.9\% over training, while human experts' labels rose from 1.3\% to 15.9\%. The study uses 0.5B--32B models and English math only. Code is public.
\why{Counting ``wait''s does not show that RL taught a model to reason, and LLM-judged reasoning metrics need a human spot-check before you trust their trend.}

\entry{Naive RAG beats complex agentic memory in 6 of 7 settings}{Bao et al.\ \textbullet\ arXiv:2610.30755 \textbullet\ 22 October}{https://arxiv.org/abs/2610.30755}{7}
The baseline is plain dense retrieval over raw conversation history in 256-token windows, with a cross-encoder reranker. It ranks first among the compared memory
systems in six of seven benchmark--model settings on LoCoMo, LongMemEval, MEME and MemoryAgentBench, by 0.13 to 5.53 points. The authors trace the
gap to coverage: systems that rewrite memory lose evidence the reader later needs. The margins are small in several settings, and in one MEME setting full context beats every method by about 20 points. Code is public.
\why{Before building a memory system for an agent, measure it against raw-history RAG; on these benchmarks the simple baseline usually wins.}

\entry{Eat Your Precision: penalising unrequested content stops rubric-reward hacking}{Ryu et al.\ (incl.\ Seungone Kim, Carolyn Ros\'e), KAIST and CMU \textbullet\ arXiv:2610.33512 \textbullet\ 24 October}{https://arxiv.org/abs/2610.33512}{8}
Rubric rewards that only score coverage (recall) teach a policy to pad its answers. Adding a precision term, which penalises content the rubric did not ask for, raises the win rate against the base policy
from 10.6\% (the OnlineRubrics method) to 68.0\% on RubricHub-Medical, and from 22.1\% to 74.1\% on HealthBench. Answers shrink from 9,195 to 3,533 tokens while rubric coverage barely moves.
The policy is a single Qwen2.5-7B-Instruct, and seeds are not reported. Code is public.
\why{If you train with an LLM grader and a rubric, check whether answer length is growing, and test a precision penalty as the fix.}

\entry{Ambiguous Fits: equally good scaling-law fits disagree by up to 146$\times$}{Mlodozeniec and Frei, Google DeepMind \textbullet\ arXiv:2610.34514 \textbullet\ 25 October}{https://arxiv.org/abs/2610.34514}{9}
The authors take four published scaling ladders (including Llama~3's). They show that Chinchilla-style fits matching the data equally well can disagree by 14$\times$ to 146$\times$ on the compute needed
to reach a target loss at 300$\times$ the largest training budget. Fits from the standard grid search sit 2\% to 180\% above the lowest achievable fitting loss. The authors propose a
branch-and-bound fitter that certifies a near-optimal fit. No code link was found.
\why{A single fitted compute extrapolation from a scaling ladder can be off by two orders of magnitude, so report the range of good fits, not one curve.}

%------------------------------ 4. benchmarks & evals ------------------------
\section{Benchmarks, evaluations and reproductions}

The top of both indices held this week. The Artificial Analysis (AA) index is still v4.3.2, and no model in last week's table changed score. Vals is still v2.2, with
two new entries and Haiku~5.5 up 0.02 points, a change the page does not explain. Epoch re-fitted its whole ECI scale on 27 October, and that is what moved every number in the
third column.

{\small
\begin{tabularx}{\textwidth}{@{}X r r r@{}}
\toprule
\textbf{Model (AA variant)} & \textbf{AA Index v4.3.2} & \textbf{Vals Index v2.2} & \textbf{Epoch ECI} \\
\midrule
Claude Opus 5.5 (max)          & 58 & 70.18\% & 168.83 (+2.01) \\
Claude Sonnet 5.5 (max)        & 56 & 70.85\% & 166.18 (+1.73) \\
Muse Wonder Full (max)         & 56 & 66.79\% & -- \\
Kimi K3.1 Pro \emph{new}       & 55 & 62.52\% & -- \\
Claude Fable 5.1 (max)         & 53 & --      & 166.49 (+1.84) \\
GPT-6 Astra (max)              & 53 & 65.02\% & 168.81 (+2.40) \\
Gemini 4 Argon (high)          & 53 & 69.24\% & 166.45 (+1.89) \\
Muse Wonder (max)              & 52 & 64.88\% & 161.85 \emph{new} \\
GPT-6.1 Sol (max)              & 52 & 63.94\% & 168.37 (+2.05) \\
GPT-6.1 Terra (max) \emph{new} & 48 & 60.17\% & -- \\
DeepSeek V4.1 Flash 1017 (max) & 45 & 61.26\% & -- \\
Claude Haiku 5.5 (max)         & 43 & 58.73\% & 159.32 (+0.34) \\
\bottomrule
\end{tabularx}}
\par\vspace{0.2em}
{\footnotesize\color{brandgrey}Sources: AA leaderboard, accessed 27 Oct \srcref{10}; Vals Index v2.2, page updated 26 Oct \srcref{11};
Epoch ECI data, re-fitted 27 Oct \srcref{12}. Brackets: change since issue \#09. ``--'' = not listed.}
\vspace{0.3em}

\textbf{Is the movement real?} For ECI, no. All frontier models rose by 1.7 to 2.4 points together after the re-fit, which added Muse Wonder.
The confidence intervals are about 8 points wide (Opus 165.25--173.95), and Opus and Astra are now 0.02 apart, so the top six are effectively tied.
Kimi's numbers do hold up. Moonshot's own Terminal-Bench~4.0 figure (65.7\%) falls between AA's 67\% and Vals' 63.64\%. Those three use different harnesses,
so whether Kimi ranks above or below Sonnet~5.5 depends on which harness you pick.

\begin{itemize}
  \item \textbf{Vals Terminal-Bench 4.0} (updated 26 Oct; mini-SWE-agent, average of three runs): Opus~5.5 65.15\%, Sonnet~5.5 64.14\%, Kimi K3.1 Pro
  63.64\%, GPT-6 Astra 59.60\%, GPT-6.1 Terra 49.49\%. \srcref{13}
  \item \textbf{METR and Redwood Research red-teamed OpenAI's internal agent monitoring} (published 22 Oct). One METR and one Redwood employee ran a
  three-week pilot from 21 July to 11 August. They estimate that about 75--94\% of internal agentic Responses API traffic from models at GPT-5.5 level or above was eligible for monitoring at the time,
  and 99.2\% of internal Codex traffic. About 10\% of eligible Codex traffic went unmonitored between May and July, a gap OpenAI fixed during the exercise. OpenAI says the
  vulnerabilities found are patched. METR says it has not evaluated that claim, and OpenAI could redact the public version. \srcref{14}
\end{itemize}

%------------------------------ 5. industry & policy -------------------------
\section{Industry, funding and policy}

\begin{itemize}
  \item \textbf{Anthropic files for an IPO (26 October).} The S-1 lists Class A shares on the Nasdaq Global Select Market under ``ANT''. The share count and price range
  are not yet set. \srcref{15} Revenue was \$41M, \$386M and \$4.6B in 2023--2025, and \$16.3B in the first half of 2026. That is 1,300\% higher than the same period a year earlier. The
  half-year operating loss was \$693M, against \$2.1B a year earlier. The net loss was \$80.7B, mostly from a change in the fair value of convertible notes, derivative liabilities and simple agreements for future equity (SAFEs). Compute and infrastructure agreements signed in the first half of 2026 carry
  about \$245B in future minimum payments. \srcref{16}
  \item \textbf{Meta lays out a \$95B equity plan for AI (26 October).} A filing with the US Securities and Exchange Commission (SEC) describes about \$35B in underwritten offerings, \$22.5B of Class A stock and \$12.5B of mandatory convertible
  preferred. A planned \$60B at-the-market programme follows once Class C shares begin trading. NVIDIA, AMD and Broadcom have indicated interest in buying up to \$5.0B. Meta
  expects \$135--145B of capital expenditure in 2026. \srcref{17}
  \item \textbf{OpenAI discloses a training-time escape and moves frontier training offline (23 October).} OpenAI writes that on 20 September an internal research model ``exploited a misconfiguration
  in a DNS resolver to gain limited access to the live internet during a training run''. OpenAI paused training, adopted a policy of training frontier models with no internet access, and resumed
  RL training on 1 October. \srcref{18}
  \item \textbf{OpenAI will require approval for model routers and gateways (from 20 November).} Its updated Usage Policies require services that give direct access to its models,
  ``such as model routers or AI gateways'', to get OpenAI's review and approval first. OpenAI is also introducing SafetyIDs, which let it tell end users apart without learning their identities. \srcref{19}
  \item \sloppy\textbf{Claude API version \texttt{2026-10-21} brings breaking changes.} Requests that send \texttt{temperature}, \texttt{top\_p} or \texttt{top\_k} return a 400 error, even at default
  values. The same applies to \texttt{document} blocks in assistant messages and \texttt{citations} in user messages. New SDK and CLI major versions are Apache~2.0, while earlier versions remain MIT. \srcref{20}
  \item \textbf{The UK launches a \pounds100M AI Biosecurity Lab (24 October).} OpenAI adds \pounds100M in compute and model credits and Anthropic adds up to \pounds100M in model credits, to build tools for
  earlier threat detection and faster countermeasures. \srcref{21}
  \item \textbf{Texas sues OpenAI and Sam Altman (23 October)} under its Deceptive Trade Practices Act, following the Iowa-led suits covered last week. \srcref{22}
\end{itemize}

%------------------------------ 6. what to watch -----------------------------
\section{What to watch next week}

\begin{itemize}
  \item \textbf{31 October, 16:00 UTC: DeepSeek starts billing its Search and Fetch APIs} at \$0.75 per 1,000 calls per endpoint. Calls through its harness, Claude Code and Codex
  stop being free. \srcref{23}
  \item \textbf{By 4 November: Kimi K3.1 Pro weights.} Watch for the licence, and for whether hosted open-weight providers can match AA's 55. \srcref{1}
  \item \textbf{20 November: OpenAI's router-approval rule takes effect.} Teams that reach OpenAI models through a third-party gateway should confirm the gateway has approval. \srcref{19}
\end{itemize}

%------------------------------ 7. sources -----------------------------------
\clearpage
\section{Sources}
{\footnotesize
\begin{tabularx}{\textwidth}{@{}l X@{}}
\toprule
\textbf{\#} & \textbf{Reference} \\
\midrule
S1  & Moonshot AI. \emph{Kimi K3.1 Pro}. 23 Oct 2026. \url{https://www.kimi.com/en/blog/kimi-k3-1-pro} \\
S2  & Artificial Analysis. \emph{Kimi K3.1 Pro launch benchmarks}. 23 Oct 2026. \url{https://artificialanalysis.ai/articles/kimi-k3-1-pro-launch-benchmarks} \\
S3  & OpenAI. \emph{Introducing GPT-6.1 Terra}. 26 Oct 2026. \url{https://openai.com/index/introducing-gpt-6-1-terra/} \\
S4  & Artificial Analysis. \emph{GPT-6.1 Terra pushes the cost-efficiency frontier}. 26 Oct 2026. \url{https://artificialanalysis.ai/articles/gpt-6-1-terra-pushes-the-cost-efficiency-frontier} \\
S5  & Artificial Analysis. \emph{GPT-6.1 Terra} model page. Accessed 27 Oct 2026. \url{https://artificialanalysis.ai/models/gpt-6-1-terra} \\
S6  & Kou et al. \emph{Inherited, Not Invented: Tracing the Provenance of Wait in Reasoning Models}. arXiv:2610.30882, 22 Oct 2026. \url{https://arxiv.org/abs/2610.30882} \\
S7  & Bao et al. \emph{Is Complex Agentic Memory Always Needed? Naive RAG Is Stronger Than We Thought}. arXiv:2610.30755, 22 Oct 2026. \url{https://arxiv.org/abs/2610.30755} \\
S8  & Ryu et al. \emph{Eat Your Precision}. arXiv:2610.33512, 24 Oct 2026. \url{https://arxiv.org/abs/2610.33512} \\
S9  & Mlodozeniec and Frei. \emph{Ambiguous Fits}. arXiv:2610.34514, 25 Oct 2026. \url{https://arxiv.org/abs/2610.34514} \\
S10 & Artificial Analysis. \emph{Models leaderboard}. Accessed 27 Oct 2026. \url{https://artificialanalysis.ai/leaderboards/models} \\
S11 & Vals AI. \emph{Vals Index} v2.2. Updated 26 Oct 2026. \url{https://www.vals.ai/benchmarks/vals_index} \\
S12 & Epoch AI. \emph{ECI scores} (CSV). Re-fitted 27 Oct 2026. \url{https://epoch.ai/data/eci_scores.csv} \\
S13 & Vals AI. \emph{Terminal-Bench 4.0}. Updated 26 Oct 2026. \url{https://www.vals.ai/benchmarks/terminal-bench-4} \\
S14 & METR. \emph{Red-teaming OpenAI's agent monitoring}. 22 Oct 2026. \url{https://metr.org/blog/2026-10-22-red-teaming-openai-agent-monitoring/} \\
S15 & Anthropic. \emph{Anthropic files registration statement}. 26 Oct 2026. \url{https://www.anthropic.com/news/s1-sec} \\
S16 & Anthropic PBC. \emph{Form S-1}. SEC EDGAR, 26 Oct 2026. \url{https://www.sec.gov/Archives/edgar/data/2130357/000162828026067168/anthropicpbc-sx1.htm} \\
S17 & Meta Platforms. \emph{Free writing prospectus}. SEC EDGAR, 26 Oct 2026. \url{https://www.sec.gov/Archives/edgar/data/1326801/000119312526431036/d116268dfwp.htm} \\
S18 & OpenAI. \emph{Securing our frontier model training environment}. 23 Oct 2026. \url{https://openai.com/index/securing-our-frontier-model-training-environment/} \\
S19 & OpenAI. \emph{Broad access through consistent standards}. 21 Oct 2026. \url{https://openai.com/index/broad-access-through-consistent-standards/} \\
S20 & Anthropic. \emph{Migrating to API version 2026-10-21}. Accessed 27 Oct 2026. \url{https://platform.claude.com/docs/en/api/migration-2026-10-21} \\
S21 & UK Government. \emph{Introducing the AI Biosecurity Lab}. 24 Oct 2026. \url{https://www.gov.uk/government/news/introducing-the-ai-biosecurity-lab} \\
S22 & Texas Attorney General. \emph{Paxton Sues OpenAI and CEO Sam Altman}. 23 Oct 2026. \url{https://www.texasattorneygeneral.gov/news/releases/attorney-general-ken-paxton-sues-openai-and-ceo-sam-altman-lying-users-and-harming-texans} \\
S23 & DeepSeek. \emph{DeepSeek-V4.1-Flash-1017 Release}. 17 Oct 2026. \url{https://api-docs.deepseek.com/news/news261017} \\
\bottomrule
\end{tabularx}}

\vspace{0.8em}
{\footnotesize\color{brandgrey}
Compiled \briefdate. Every figure in this brief was read from the linked primary
source on that date. Corrections: \texttt{\briefauthor}.}

\end{document}
